% ============================================================
% INTRODUCTION GÉNÉRALE
% ============================================================
\chapter*{Introduction générale}
\addcontentsline{toc}{chapter}{Introduction générale}

\section*{Contexte}

Au Bénin, la Contraction de texte constitue l'une des épreuves les plus redoutées du Baccalauréat (séries A, B, C, D). Elle exige de l'élève une triple maîtrise : comprendre fidèlement les idées de l'auteur, les reformuler sans plagiat, et respecter un quota de mots strict — le quart du texte original, avec une tolérance de \(\pm10\%\). Cette épreuve est systématiquement sous-préparée faute de moyens de correction individualisée dans des classes de plus de 60 élèves.

L'élève reçoit sa copie annotée des semaines après la remise, brisant la boucle d'apprentissage cognitive essentielle à la progression. Aucun outil numérique existant ne propose un entraînement structuré à la Contraction de texte en s'appuyant sur une grille d'évaluation académique pour préparer l'apprenant au Baccalauréat.

\section*{Problématique}

Comment concevoir et évaluer expérimentalement un pipeline NLP hybride — combinant traitement déterministe local et modèle de langage large — capable d'évaluer automatiquement la Contraction de texte en s'appuyant sur une grille d'entraînement académique, de générer un feedback formatif personnalisé, et de modéliser la progression adaptative de l'apprenant en vue de sa préparation au Baccalauréat ?

Trois sous-questions structurent cette problématique :

\begin{enumerate}
    \item Quelle combinaison de métriques de similarité (LCS, Cosinus TF-IDF) offre la meilleure performance pour la détection du plagiat et la validation du quota dans le contexte de la Contraction de texte en français ?
    \item Comment concevoir des prompts Few-Shot efficaces pour qu'un LLM génère un feedback formatif pertinent et fidèle à la grille d'entraînement académique, sans hallucinations sur la forme ?
    \item Comment modéliser dynamiquement la progression d'un apprenant en Contraction de texte et adapter la difficulté en conséquence ?
\end{enumerate}

\section*{Objectifs}

\subsection*{Objectif général}

Concevoir, implémenter et évaluer expérimentalement \AZIZA, un pipeline NLP hybride d'évaluation automatique de la Contraction de texte, avec une démonstration fonctionnelle via une interface Flask.

\subsection*{Objectifs spécifiques}

\begin{enumerate}
    \item \textbf{Pipeline NLP — Forme :} décompte strict des mots et calcul de similarité textuelle (LCS, Cosinus) pour la détection du plagiat.
    \item \textbf{Expérimentation comparative :} évaluer les performances des métriques (Précision, Rappel, F1-score, corrélation de Spearman).
    \item \textbf{Pipeline LLM — Fond :} concevoir des prompts Few-Shot avec chaîne de pensée (CoT) pour évaluer la fidélité sémantique et générer un feedback formatif structuré.
    \item \textbf{Évaluation du feedback LLM :} mesurer la pertinence pédagogique via un annotateur expert (Cohen's Kappa).
    \item \textbf{Modélisation de l'apprenant :} machine à états à 3 paliers (Découverte, Consolidation, Maîtrise) avec un algorithme de progression adaptatif.
    \item \textbf{Interface de démonstration :} application Flask avec API REST, Swagger, et déploiement PythonAnywhere.
\end{enumerate}

\section*{Contributions}

Ce mémoire apporte les contributions suivantes :

\begin{itemize}
    \item Un \textbf{pipeline NLP hybride original} combinant métriques déterministes (Forme /7) et LLM (Fond /3), garantissant l'absence d'hallucination sur les critères quantitatifs.
    \item Une \textbf{expérimentation scientifique rigoureuse} sur 30 contractions annotées, avec des métriques d'évaluation complètes (Pearson, Kappa, F1-score, courbe ROC).
    \item Un \textbf{corpus annoté} de 30 contractions réparties en 3 niveaux de difficulté (correct, partiellement correct, incorrect), issues des annales du Baccalauréat.
    \item Un \textbf{système de feedback formatif} généré par LLM, structuré en trois volets : diagnostic de l'erreur, micro-cours pédagogique, et exemple de reformulation.
    \item Une \textbf{architecture logicielle modulaire} (API REST, rate limiting, Swagger) déployée sur PythonAnywhere, prête pour l'intégration dans une application web complète.
\end{itemize}

\section*{Organisation du mémoire}

Ce mémoire est structuré en cinq chapitres. Le \textbf{Chapitre 1} dresse l'état de l'art sur l'évaluation automatique de textes, les systèmes de tutorat intelligent et le TALN francophone. Le \textbf{Chapitre 2} présente la méthodologie et l'architecture du pipeline hybride. Le \textbf{Chapitre 3} détaille l'implémentation, des décisions techniques aux déploiement. Le \textbf{Chapitre 4} expose les expérimentations et leurs résultats. Le \textbf{Chapitre 5} présente les tests et la validation. Une conclusion générale synthétise les résultats et ouvre sur les perspectives d'évolution.
